%% paper.tex
%% AIDI 2026 Workshop @ FLICS 2026 — IEEE conference template
%% Compiled with: pdflatex paper.tex; bibtex paper; pdflatex paper.tex (x2)

\documentclass[conference]{IEEEtran}
\IEEEoverridecommandlockouts
\usepackage{cite}
\usepackage{amsmath,amssymb,amsfonts}
\usepackage{graphicx}
\usepackage{textcomp}
\usepackage{xcolor}
\usepackage{booktabs}
\usepackage{multirow}
\usepackage{array}
\usepackage{url}
\usepackage{siunitx}

% Convenience macros
\newcommand{\dlidar}{\ensuremath{\Delta_{\mathrm{LiDAR}}}}
\newcommand{\dlidarw}{\ensuremath{\Delta_{\mathrm{LiDAR}}^{\mathrm{wi}}}}
\newcommand{\dB}{\,\mathrm{dB}}
\newcommand{\dBm}{\,\mathrm{dBm}}

\begin{document}

\title{LiDAR-Augmented WiFi Signal Strength Prediction:\\A Use Case Study on Industrial AGVs}

\author{%
\IEEEauthorblockN{Myroslav Mishchuk\textsuperscript{1,2},
                  Olena Pavliuk\textsuperscript{1,2},
                  Ireneusz Smo{\l}ka\textsuperscript{1},
                  Vasyl Tesliuk\textsuperscript{2},
                  Rafa{\l} Cupek\textsuperscript{1}}
\IEEEauthorblockA{\textsuperscript{1}Silesian University of Technology, Poland\\
\textsuperscript{2}Lviv Polytechnic National University, Ukraine}
}

\maketitle

\begin{abstract}
Industrial automated guided vehicles (AGVs) require predictable WiFi
connectivity for safety telemetry and fleet coordination, but signal
strength in industrial workspaces is hard to predict from access point
(AP) distance alone due to waveguiding, blockage, and multipath. Since
industrial AGVs are routinely equipped with safety LiDARs that capture
the local environmental geometry, the natural hypothesis is that
LiDAR-derived features can augment WiFi signal strength prediction
beyond what AP-relative geometry alone provides. This paper provides
a use case study on data gathered from a deployed industrial AGV
(three routes, 681,593 matched samples) using cross-session and
within-session generalization protocols with five robustness
controls: an XGBoost
hyperparameter sweep, a LightGBM framework cross-check, a 1-meter
spatial buffer-zone exclusion, a depth-constrained over-fitting check,
and LiDAR-feature permutation tests. Across all configurations,
ego-frame LiDAR features did not improve on a 15-feature baseline of
position, telemetry, and AP-relative geometry: the cleanest LiDAR-free
model achieves an RMSE of 1.99 dB on signal\_power, approximately 3 dB
below the dataset's noise floor. The paper provides the first
controlled empirical study of LiDAR-augmented WiFi signal strength
prediction on industrial AGVs and offers practical deployment guidance
for fleet operators.
\end{abstract}

\begin{IEEEkeywords}
Industrial AGV, WiFi link quality, LiDAR, signal-power prediction,
XGBoost, SHAP, explainable AI (XAI)
\end{IEEEkeywords}

%==========================================================================
\section{Introduction}\label{sec:intro}
%==========================================================================

With the ongoing adoption of Industry 4.0 and 5.0 paradigms,
automated guided vehicles (AGVs) are increasingly deployed in
warehouse, manufacturing, and distribution
environments~\cite{fragapane2021planning}. Reliable wireless
communication supports the coordination functions of an AGV fleet,
including protective telemetry, traffic management, mission updates,
and remote intervention; unpredictable signal quality can translate
into idle time, near-collisions, or unplanned safety stops. WiFi
remains the dominant choice for these links because of its cost,
infrastructure availability, and coverage characteristics, but the
conditions under which industrial WiFi performs reliably differ from
those of office or residential
deployments~\cite{cena2025wifi,schmieder2023agvdatasets}. Metallic
shelving, structural columns, loading equipment, and moving forklifts
together produce a propagation environment dominated by waveguiding,
blockage, and multipath. Empirical log-distance models, even with
precisely surveyed access point (AP) coordinates, fit indoor
industrial measurements with considerably more residual variance than
outdoor or campus deployments~\cite{cena2025wifi,formis2025wifi}.
Field-deployable predictors of received signal power along a planned
AGV trajectory are therefore operationally valuable: they let a fleet
manager pre-compute coverage, flag risk zones, and schedule jobs
around predicted dropout regions~\cite{gulia2025evowisva}.

Industrial AGVs designed to comply with ISO 3691-4~\cite{iso3691-4}
are typically equipped with 2D safety laser scanners (commonly
referred to as 2D safety LiDARs) for personnel detection and
protective-field monitoring; in many platforms the same scanners
also feed the on-board localization stack. Although their primary
function is safety, these scanners continuously produce a wide-angle
range scan of the AGV's local environment, capturing corridor
geometry, blockage along arbitrary directions, and clutter density.
Indoor WiFi propagation is shaped by exactly these geometric
properties, which motivates the question of whether LiDAR-derived
environmental features can augment a signal-strength prediction
model beyond what AP-relative geometry alone provides --- and whether
the added information justifies the integration effort. Recent work
has shown that point-cloud-derived environmental features can
improve mmWave-band channel
prediction~\cite{ohta2023pointcloud,nist2024mmwave,zhang2025mmlscm,
clark2022propeml}, and SHAP-based attribution has been combined with
gradient-boosted regression for explainable cellular path-loss
models~\cite{xu2025xgboost5g,kiouvrekis2025explainable,
balogun2022satelliteshap}. Existing AGV WiFi link-quality
predictors~\cite{formis2025wifi,formis2024deep,colletto2023ann} use
telemetry alone and do not incorporate environmental geometry. To
our knowledge, the intersection of LiDAR-derived environmental
features, WiFi link-quality prediction, SHAP-based explanation, and
industrial-AGV deployment is not yet explored
(Section~\ref{sec:rw}).

This paper presents a use case study, conducted on data gathered
from a deployed industrial AGV across three routes, that examines
whether ego-frame LiDAR features augment such a baseline. The
dataset, source code, analysis scripts, and result reports are
released to support reproduction and follow-up
work\footnote{Dataset, source code, analysis scripts, and reports:
\mbox{\url{https://github.com/Myroslav437/FLICS2026-LiDAR-Derived-WiFi-Quality-Prediction}}.}.

The paper addresses the following research questions (RQ):
\begin{itemize}
\setlength\itemsep{2pt}
\item \textbf{RQ1.} Do LiDAR-derived environment features improve
  received signal-power prediction over a baseline of AGV position,
  telemetry, and AP-relative geometry, in cross-session
  generalization?
\item \textbf{RQ2.} Do LiDAR-derived environment features improve
  prediction in within-session leave-region-out generalization, after
  controls for spatial autocorrelation and tree-depth over-fitting?
\item \textbf{RQ3.} Are the observed effects robust to feature-group
  substitution, hyperparameter and framework choice, and to a
  LiDAR-feature permutation control?
\end{itemize}

The contributions of this paper are summarized as follows:
\begin{itemize}
\setlength\itemsep{2pt}
\item The first controlled empirical evaluation of LiDAR-augmented
  WiFi signal strength prediction on data gathered from a deployed
  industrial AGV is performed.
\item A feature-group ablation procedure is developed that separates
  the case where a feature group provides no useful information from
  the case where it is substitutable for another retained group.
\item Practical deployment guidance is provided for industrial AGV
  operators on the minimal feature set required for accurate
  within-session signal-power prediction in the studied scenario.
\end{itemize}

The remainder of the paper is organized as follows.
Section~\ref{sec:rw} reviews related work. Section~\ref{sec:data}
describes the dataset, its preparation, and the feature stack.
Section~\ref{sec:method} presents the experimental methodology.
Sections~\ref{sec:cross-session} and~\ref{sec:within-session} report
the cross-session and within-session experiments, respectively.
Section~\ref{sec:shap} presents the feature attribution analysis.
Section~\ref{sec:discussion} discusses the implications and
limitations, and Section~\ref{sec:conclusion} concludes.

%==========================================================================
\section{Related Work}\label{sec:rw}
%==========================================================================

A coherent body of work on industrial WiFi link-quality prediction
has used telemetry-only features such as recent received-power
history, throughput, and protocol counters, without explicit
modelling of the surrounding workspace. Linear combinations of
exponential moving averages~\cite{formis2023moving} and shallow
neural networks~\cite{colletto2023ann,scanzio2022ann} reach useful
accuracy on warehouse-deployed equipment, and recent extensions
replace the shallow models with deep
architectures~\cite{formis2024deep,formis2025wifi}; the feature
inventory across these works remains essentially telemetric. Broader
surveys of machine learning for IEEE
802.11~\cite{wilhelmi2024aiml,szott2022wifimeetsml} catalogue similar
feature sets and identify environmental modelling as an
under-investigated direction, particularly in industrial deployments
where waveguiding and blockage dominate; the fundamental challenges
of WiFi reliability for moving equipment in such environments are
also documented in~\cite{cena2025wifi}. The cumulative picture is
that telemetry features alone leave residual variance that the prior
art has not been able to explain away, which motivates the search
for additional features that capture environmental structure.

Recent work has explored that direction predominantly in the
millimetre-wave (mmWave) band rather than indoor WiFi. Raw point
clouds have been used to anticipate line-of-sight blockage in mmWave
networks~\cite{ohta2023pointcloud}; a NIST measurement-based
predictor fuses LiDAR scans with panoramic imagery in a deep
regressor for 60\,GHz channel parameters~\cite{nist2024mmwave}; a
multi-modal neural radio radiance field combines reference signal
received power (RSRP) with LiDAR~\cite{zhang2025mmlscm}; PropEM-L
encodes raw point clouds for communication-aware exploration in
outdoor multi-robot systems~\cite{clark2022propeml}; and
GPU-accelerated point-cloud ray tracing has been demonstrated for
outdoor channel prediction~\cite{guo2021gpuraytrace}. The reported
gains are substantial, but the mmWave regime differs materially from
the 2.4/5\,GHz indoor industrial case of interest here: mmWave is
dominated by direct-path attenuation and discrete blockage events,
whereas industrial WiFi is shaped by waveguiding through corridors
and constructive/destructive multipath at length scales much larger
than typical LiDAR feature scales. Whether geometric features that
succeed in mmWave transfer to industrial WiFi is therefore an open
empirical question.

A separate line of work has applied explainable AI --- in particular,
SHAP-based attribution --- to gradient-boosted path-loss models.
XGBoost with SHAP has been used for 5G urban path-loss prediction
with multi-modal inputs~\cite{xu2025xgboost5g};
the prediction accuracy and explanation power of SHAP-based and
additive models have been compared in a university-campus
environment~\cite{kiouvrekis2025explainable}; and hand-crafted versus
convolutional features extracted from satellite imagery have been
compared for outdoor path-loss
prediction~\cite{balogun2022satelliteshap}. Across these works the
regression target is outdoor cellular or campus WiFi, and the
environmental modality is imagery rather than LiDAR; the
methodology is transferable, but the transfer to indoor industrial
AGV deployments has not, to the best of our knowledge, been
carried out.

Dedicated LiDAR/WiFi sensor-fusion work in the AGV domain has so far
targeted localization, simultaneous localization and mapping, and
navigation drift compensation~\cite{li2025lidaragv,
nguyen2022wifilidar}, not link-quality prediction. This line of work
confirms that the two sensor streams are routinely available
together on production AGVs --- which makes the augmented prediction
investigated here practically attractive --- but does not address
the prediction question.

To the best of our knowledge, the intersection of LiDAR-derived
environmental features, WiFi link-quality prediction, SHAP-based
explanation, and industrial-AGV deployment has not yet been studied
empirically. The use case study presented in this paper is intended
to contribute to that gap with the experimental controls necessary
to support a credible deployment-relevant conclusion.

%==========================================================================
\section{Dataset Collection and Preparation}\label{sec:data}
%==========================================================================

This section describes the experimental platform, the data-gathering
campaign, the pre-processing steps required to convert raw streams
into a training-ready dataset, and the feature stack supplied to the
predictive models.

%-------------------------------------------------------------------- 3.1
\subsection{Platform and data-gathering campaign}
The experimental platform is a single industrial AGV equipped with a
front-facing Leuze RSL\,400 270$^{\circ}$ 2D safety laser scanner
(operated at 0.2$^{\circ}$ angular resolution; 1{,}350 active beams
in a 2{,}700-slot buffer) and a Navitrol on-board navigation and
supervision system (NNS) by Navitec Systems. Sample rates of the
LiDAR, telemetry, and NNS streams are summarized in
Table~\ref{tab:sessions}. The set of features extracted from the raw
streams and supplied to the predictive models is described in
Section~\ref{sec:features} and listed in
Table~\ref{tab:features}. The facility is served by a single WiFi
access point (AP) whose coordinates are surveyed in the AGV's
reference frame to within several centimetres.

Three data-gathering sessions were collected on a deployed AGV. The
sessions cover three distinct routes, summarized in
Table~\ref{tab:sessions} and visualized in
Fig.~\ref{fig:trajectories}. S1 and S2 share the route's reference
frame; S3 was recorded on a different route in a different sub-area
of the facility. S1 is the largest session, both in duration and in
spatial coverage, and is used for the within-session experiments.

\begin{table}[!t]
\caption{Three-session dataset overview.}
\label{tab:sessions}
\centering
\renewcommand{\arraystretch}{1.0}
\setlength{\tabcolsep}{4pt}
\footnotesize
\begin{tabular}{l c c c c c}
\toprule
Session & $n_{\text{pairs}}$ & dur.\ (h) & $\overline{P}$ (dBm) & $\sigma_P$ (dB) & motion \% \\
\midrule
S1    & 280{,}025 & 3.19 & $-36.91$ & 10.52 & 37.6 \\
S2    & 168{,}940 & 1.93 & $-39.66$ & 7.35  & 39.3 \\
S3    & 232{,}628 & 2.75 & $-31.07$ & 6.22  & 61.3 \\
\midrule
total & 681{,}593 & 7.87 & --- & --- & --- \\
\bottomrule
\end{tabular}
\\[3pt]
\raggedright
\footnotesize
LiDAR scans at $\sim$25\,Hz; telemetry at 31\,Hz;
NNS pose at 4.4\,Hz. ``Motion \%'' is the fraction of recorded time
during which the AGV is motion-active ($|\texttt{speed\_mps}| >
0.1\,\mathrm{m/s}$).
\vspace{-0.6em}
\end{table}

% TODO: figure showing AGV trajectory coloured by signal_power, three sessions side by side
\begin{figure}[!t]
\centering
% \includegraphics[width=\columnwidth]{figs/trajectories_signalpower.pdf}
\fbox{\parbox[c][2.6cm][c]{0.92\columnwidth}{\centering\itshape
[Placeholder Fig.~\ref{fig:trajectories}: AGV trajectory in the
session reference frame, coloured by \texttt{signal\_power}, for the
three sessions side-by-side. Hot/cold zones are visibly persistent
within each session.]}}
\caption{AGV trajectory coloured by \texttt{signal\_power}.}
\label{fig:trajectories}
\vspace{-1.0em}
\end{figure}

%-------------------------------------------------------------------- 3.2
\subsection{Time synchronization}
LiDAR scans and telemetry samples originate from independent
hardware and software stacks with their own clocks. Even moderate
temporal drift between the two streams biases any subsequent
learning, because LiDAR-derived features would describe a slightly
different position from the one the telemetry suggests. Time
alignment is therefore a prerequisite for the use case study and is
performed before any feature engineering.

A per-session offset $\hat{\tau}$ is estimated as the value that
maximizes the cross-correlation between a telemetry signature and a
LiDAR-side scan-to-scan dissimilarity signal, computed on a common
50\,Hz grid with parabolic peak refinement. The estimation is
performed once per session, and the resulting confidence interval is
quantified via moving-block bootstrap. The hypothesis that a single
offset applies to all three sessions is then tested against the
per-session estimates with Cochran's $Q$ test ($Q = 8.06$, $df = 2$,
$p = 0.018$); the test rejects the equal-offsets hypothesis, so
per-session offsets are used in the rest of the pipeline. The
estimated offsets, applied as constant per-session shifts to align
the two streams, are
$\hat{\tau}_{\mathrm{S}1} = +0.250$\,s,
$\hat{\tau}_{\mathrm{S}2} = +0.432$\,s, and
$\hat{\tau}_{\mathrm{S}3} = +0.879$\,s. After alignment and joining,
97.0\% of LiDAR scans (681{,}593 of 702{,}624) have a successfully
matched telemetry sample and are retained in the training-ready
dataset. Fig.~\ref{fig:timesync} shows the cross-correlation curve
for one session and the estimated peak.

% TODO: figure showing time-sync result on one session
\begin{figure}[!t]
\centering
% \includegraphics[width=\columnwidth]{figs/timesync.pdf}
\fbox{\parbox[c][2.4cm][c]{0.92\columnwidth}{\centering\itshape
[Placeholder Fig.~\ref{fig:timesync}: cross-correlation between the
telemetry signature and the LiDAR-side scan-to-scan dissimilarity
signal, with the bootstrap-estimated peak marked, on one session.]}}
\caption{Time synchronization between LiDAR and telemetry streams.}
\label{fig:timesync}
\vspace{-1.0em}
\end{figure}

%-------------------------------------------------------------------- 3.3
\subsection{Anomaly handling}\label{sec:anomaly}
The downstream models use the AGV's reported position $(x_m, y_m)$
to compute every AP-relative feature. When the on-board localization
is struggling --- typically during initialization, after a manual
reposition of the AGV, in regions with feature-poor environments, or
under sensor degradation --- the reported position becomes
unreliable, and so do the features derived from it. Such rows must
be excluded so that the feature derivation remains physically
meaningful.

The NNS exposes a position-confidence score
\texttt{nns\_position\_confidence} on a 0--100 scale, with a high
value indicating that NNS estimates its own pose to be reliable. The
distribution of this score across the entire dataset is sharply
bimodal (Fig.~\ref{fig:confidence-hist}): a narrow high-confidence
mode at 95--100, where the localization is trustworthy, and a smaller
low-confidence mode at 19--23, where it is not. A 5-bin moving average
of the histogram identifies a clear saddle between the two modes at
$T^{\star} = 35$. Rows with confidence below this threshold are
excluded from training and evaluation. The threshold removes 74{,}429
of the 681{,}593 matched samples (10.92\%) and leaves the remaining
89.08\% for downstream use.

% TODO: figure showing the bimodal confidence histogram with T* threshold marked
\begin{figure}[!t]
\centering
% \includegraphics[width=\columnwidth]{figs/confidence_hist.pdf}
\fbox{\parbox[c][2.2cm][c]{0.92\columnwidth}{\centering\itshape
[Placeholder Fig.~\ref{fig:confidence-hist}: histogram of
\texttt{nns\_position\_confidence} across the full dataset (bimodal),
with the saddle threshold $T^{\star} = 35$ marked.]}}
\caption{Distribution of \texttt{nns\_position\_confidence} and the
exclusion threshold $T^{\star} = 35$.}
\label{fig:confidence-hist}
\vspace{-1.0em}
\end{figure}

%-------------------------------------------------------------------- 3.4
\subsection{Dataset noise floor}\label{sec:noisefloor}
Reported model accuracy values are interpretable only against a
reference for the irreducible variability of the data. We therefore
characterize the dataset noise floor at the spatial resolution at
which features are realistically deployable, namely 0.5\,m position
binning: if two visits to the same 0.5-m cell at different times
yield substantially different \texttt{signal\_power} readings, no
model that uses 0.5-m position binning alone can resolve below that
variability.

The noise-floor estimate uses the 85 cells of size 0.5\,m that contain
at least 30 rows in both S1 and S2 (the only two sessions sharing a
workspace, allowing same-cell repeat measurements). The mean per-cell
standard deviation of \texttt{signal\_power} within a session is
4.90\,dB on S1, 5.01\,dB on S2, and 4.95\,dB pooled. The median
absolute difference between the two sessions' per-cell mean values is
4.24\,dB. We therefore take 4.95\,dB as the noise floor at 0.5-m
position binning. Any predictor that resolves below this floor must be
exploiting information at finer spatial granularity than 0.5\,m ---
for instance, position at sensor resolution, telemetry, or AP-relative
geometry. We refer back to this number when interpreting model RMSE in
Sections~\ref{sec:cross-session} and~\ref{sec:within-session}.

%-------------------------------------------------------------------- 3.5
\subsection{Feature stack}\label{sec:features}
The training-ready dataset has 44 columns; 38 of them are numeric
features supplied to the predictive models, organized into five
feature groups summarized in Table~\ref{tab:features}.

\begin{table*}[!t]
\caption{Feature stack supplied to the predictive models, partitioned
into five groups (38 features in total).}
\label{tab:features}
\centering
\renewcommand{\arraystretch}{1.05}
\setlength{\tabcolsep}{5pt}
\footnotesize
\begin{tabular}{l c p{0.30\textwidth} l p{0.40\textwidth}}
\toprule
Group & \# & Feature(s) & Type / units & Description \\
\midrule
Position             & 2  & $x_m,\ y_m$                              & float / m         & AGV centroid in the workspace reference frame \\[1pt]
\midrule
\multirow{6}{*}{Telemetry} & 1 & \texttt{speed\_mps}              & float / m\,s$^{-1}$  & AGV linear speed \\
                     & 1  & \texttt{turn\_rate}                       & float / rad\,s$^{-1}$ & AGV angular velocity \\
                     & 3  & \texttt{load\_long}, \texttt{load\_mid}, \texttt{load\_short} & float / a.u.\ & strain-gauge load components on three axes \\
                     & 1  & \texttt{battery\_voltage}                 & float / V        & battery voltage \\
                     & 1  & \texttt{momentary\_current}               & float / A        & instantaneous current consumption \\
                     & 1  & \texttt{nns\_state}                       & integer code     & NNS internal state code \\[1pt]
\midrule
\multirow{4}{*}{AP-relative} & 1 & \texttt{dist\_to\_AP}            & float / m        & Euclidean distance from AGV to surveyed AP \\
                     & 2  & \texttt{sin\_AP}, \texttt{cos\_AP}        & float            & sine and cosine of the AP bearing relative to the AGV heading \\
                     & 1  & \texttt{clutter\_frac\_toward\_AP}        & float / [0,1]    & LiDAR clutter fraction in the 30$^{\circ}$ sector facing the AP (NaN if AP outside FOV) \\
                     & 1  & \texttt{is\_AP\_in\_FOV}                  & binary           & whether the AP bearing falls inside the LiDAR's empirical valid FOV \\[1pt]
\midrule
\multirow{3}{*}{LiDAR scalar} & 2 & \texttt{mean\_dist\_mm}, \texttt{dist\_p90\_mm} & float / mm   & mean and 90th-percentile beam distance over the valid 222$^{\circ}$ sector \\
                     & 2  & \texttt{clutter\_frac}, \texttt{openness\_frac} & float / [0,1] & clutter and openness ratios over the valid 222$^{\circ}$ sector \\
                     & 1  & \texttt{mean\_front\_mm}                  & float / mm       & mean beam distance in the front cone \\[1pt]
\midrule
\multirow{2}{*}{LiDAR sectoral} & 7 & \texttt{mean\_dist\_sector\_\{1..7\}\_mm} & float / mm  & mean beam distance in each of seven contiguous 30$^{\circ}$ sectors \\
                     & 7  & \texttt{clutter\_frac\_sector\_\{1..7\}}  & float / [0,1]    & clutter fraction in each of seven contiguous 30$^{\circ}$ sectors \\
\bottomrule
\end{tabular}
\vspace{-0.6em}
\end{table*}

The empirical valid LiDAR field of view is narrower than the
nameplate 270$^{\circ}$. Of the 1{,}350 active beams, 241 hit the AGV
body itself at sub-200\,mm range and are permanently masked; the
remaining 1{,}109 beams cover a contiguous 222$^{\circ}$ sector that
hosts the seven 30$^{\circ}$ sectoral aggregates listed in
Table~\ref{tab:features}. Reported metrics are stratified by
\texttt{is\_AP\_in\_FOV}: the in-FOV stratum is the regime in which
directional LiDAR features should be most useful, since the AP itself
falls within the LiDAR's valid sector and clutter measurements toward
the AP are well-defined.

%==========================================================================
\section{Methodology}\label{sec:method}
%==========================================================================

Two evaluation protocols are executed to probe complementary aspects
of generalization. The cross-session leave-one-route-out (LORO)
protocol uses three folds, each holding out one session
(Section~\ref{sec:cross-session}). The held-out session contains
trajectory patterns the model has not been trained on, including
(in fold F3) a route recorded in a workspace area that does not
appear in the training set. Because the position reference frames
do not align across the held-out route and the training routes,
cross-session models do not use the raw position $(x_m, y_m)$;
features are restricted to ego-frame LiDAR, telemetry, and
per-session AP-relative geometry.

The within-session leave-region-out (LRO) protocol uses five folds
on session S1 only --- the largest session by both spatial coverage
and sample count (Section~\ref{sec:within-session}). Regions are
obtained by $K$-means clustering on $(x_m, y_m)$ with $k = 5$.
Within-session models do use the raw position. The two protocols
together test the model's ability to transfer learned propagation
structure across operating conditions (LORO) and to extrapolate to
unseen physical regions of the same workspace (LRO).

The predictive model on both protocols is gradient-boosted regression
trees, trained with XGBoost (six hyperparameter configurations in
total, defined below) and cross-checked with LightGBM. The target
variable is the WiFi received signal power, in dBm, reported for the
matched LiDAR--telemetry sample.

%-------------------------------------------------------------------- 4.2
\subsection{Feature-group ablation}\label{sec:method-disamb}
A single ablation between a full and an ablated model is ambiguous:
no degradation can mean either (i) the removed feature group did not
help, or (ii) the group is substitutable for another retained group.
We therefore run two complementary ablations against the full model:
$M^\prime$ removes feature group $A$ (here, the LiDAR features) and
$M^{\prime\prime}$ removes feature group $B$ (here, the AP-relative
features). The pair of resulting RMSE differences is then classified,
with an absolute threshold of 0.5\,dB, into one of four verdicts:
\textit{$A$ removable}, \textit{$B$ removable}, \textit{both
removable}, or \textit{mixed/unclear}. The procedure is
feature-group-agnostic: it reports which group can be removed
without loss, independently of which group the analyst expected to
matter. Cross-session variants are denoted B5, B5$^\prime$, and
B5$^{\prime\prime}$; within-session variants are denoted W4,
W4$^\prime$, and W4$^{\prime\prime}$.

%-------------------------------------------------------------------- 4.3
\subsection{Robustness controls}\label{sec:method-controls}
The validity of the ablation-based comparison depends on several
conditions that are tested explicitly. In what follows, the primary
configuration of the XGBoost model is denoted H0 and is defined by
tree depth $6$, learning rate $\eta = 0.05$, $L_2$ regularization
$\lambda = 1.0$, $n_\text{est} = 2000$ boosting rounds with
early-stopping after $100$ non-improving rounds, and random seed
$20260427$. Validation is performed chronologically per session
(last $10\,\%$ of each training session). The robustness controls,
all evaluated on the cross-session F1 fold unless stated otherwise,
are:

\begin{enumerate}
\setlength\itemsep{1pt}
\item \textit{Hyperparameter sensitivity.} A six-configuration sweep
  of the XGBoost model: H0 (the primary configuration above);
  H1 (depth $4$, otherwise H0); H2 (slow learning rate $\eta = 0.02$,
  higher regularization $\lambda = 5$); H3 (depth $8$, otherwise H0);
  plus H1\textsubscript{rv} and H2\textsubscript{rv} with random rather
  than chronological validation splits.
\item \textit{Framework sensitivity.} A LightGBM cross-check with
  default and depth-equivalent hyperparameters, addressing whether the
  comparison is specific to XGBoost.
\item \textit{Feature-informativeness check.} A LiDAR-feature
  permutation control: a per-session block-shuffle of the 19
  ego-frame LiDAR columns that preserves their marginal and joint
  distributions but destroys row-level alignment with the target,
  evaluated both cross-session and within-session. If the
  shuffled-LiDAR model achieves the same ablation effect as the
  real-LiDAR model, the model was not extracting row-aligned
  information from the LiDAR features.
\item \textit{Single-fold robustness.} For any within-session fold
  that nominally clears the threshold: a 1-m spatial buffer-zone
  exclusion (training rows within 1\,m of any test row are dropped)
  controlling spatial autocorrelation, and a depth-4 hyperparameter
  swap controlling tree-depth over-fitting; plus their joint
  application.
\item \textit{Noise-floor reference.} The noise-floor characterization
  of Section~\ref{sec:noisefloor} is used to contextualize the
  absolute RMSE values reported in
  Sections~\ref{sec:cross-session} and~\ref{sec:within-session}.
\end{enumerate}

All training-set sizes after anomaly filtering are reported in
Sections~\ref{sec:cross-session} and~\ref{sec:within-session}.
Bootstrap 95\,\% confidence intervals (CIs) on RMSE are computed
with $B = 1000$ resamples.

%==========================================================================
\section{Cross-Session Experiment}\label{sec:cross-session}
%==========================================================================

The cross-session experiment uses three folds: F1 (S1 held out, train
on S2+S3), F2 (S2 held out, train on S1+S3), and F3 (S3 held out,
train on S1+S2; the 31\,Hz training sessions are decimated by $k = 7$
to match S3's 4.4\,Hz cadence). Per-fold post-anomaly training-set
sizes are 337{,}390 rows for F1, 403{,}876 for F2, and 50{,}232 for
F3 (the cadence match shrinks F3 to keep dynamics balanced). Six
ablation variants B0--B5 escalate from a distance-only baseline up
to the full ego-frame LiDAR + AP-relative model (Table~\ref{tab:b-variants}).
Two complementary variants implement the feature-group ablation
procedure of Section~\ref{sec:method-disamb}: B5$^\prime$ removes the
LiDAR features while keeping AP-relative geometry, and
B5$^{\prime\prime}$ removes the AP-relative features while keeping
LiDAR. F1 is treated as the primary fold throughout, because it
holds out the session with the largest spatial coverage and the
highest \texttt{signal\_power} variability
($\sigma_P = 10.52$\,dB), which together provide the most stringent
test of cross-session generalization.

\begin{table}[!t]
\caption{Cross-session ablation variants.}
\label{tab:b-variants}
\centering
\renewcommand{\arraystretch}{1.05}
\setlength{\tabcolsep}{4pt}
\footnotesize
\begin{tabular}{l c l}
\toprule
Variant & \# feat. & Feature groups \\
\midrule
B0                            & 1  & \texttt{dist\_to\_AP} only \\
B1                            & 3  & AP-relative geometry (no clutter) \\
B2                            & 8  & B1 + telemetry \\
B3                            & 22 & B2 + LiDAR scalar + \texttt{clutter\_frac\_toward\_AP} \\
B4                            & 30 & B3 + 7 sectoral mean-distance \\
B5                            & 32 & full ego-frame LiDAR + AP-relative \\
B5$^\prime$                   & 13 & B5 minus LiDAR features \\
B5$^{\prime\prime}$           & 27 & B5 minus AP-relative features \\
\bottomrule
\end{tabular}
\vspace{-0.6em}
\end{table}

%-------------------------------------------------------------------- 5.2
\subsection{Headline results}
Table~\ref{tab:cs-headline} reports the in-FOV
$\dlidar = \mathrm{RMSE}(B1) - \mathrm{RMSE}(B5)$ per fold under
the H0 configuration. The pre-registered practical-relevance
threshold of $+1\dB$ is not approached on any fold; F1 in-FOV is
decisively negative ($-0.52\dB$). The feature-group ablation
verdict (Table~\ref{tab:cs-disambig}) is LiDAR removable on every
fold: removing the 19 LiDAR features either improves the model or
hurts it by less than 0.5\,dB. AP-relative features are never
removable: B5$^{\prime\prime}$ is decisively worse than B5 on every
fold (by 0.33--2.14\,dB).

\begin{table}[!t]
\caption{Cross-session in-FOV $\dlidar$ per fold under H0.}
\label{tab:cs-headline}
\centering
\renewcommand{\arraystretch}{1.05}
\footnotesize
\begin{tabular}{l c c c c}
\toprule
Fold & $n_{\text{test}}$ & RMSE(B1) & RMSE(B5) & $\dlidar$ (95\,\% CI) \\
\midrule
F1 & 134{,}692 & 8.01 & 8.53 & $\mathbf{-0.52}$ ($[-0.53, -0.50]$) \\
F2 & 102{,}354 & 6.71 & 6.41 & $+0.30$ ($[+0.29, +0.32]$) \\
F3 & 121{,}831 & 8.87 & 7.73 & $+1.14$ ($[+1.12, +1.17]$) \\
\bottomrule
\end{tabular}
\\[3pt]
\raggedright
\footnotesize
$\dlidar > 0$ means LiDAR features improve over the AP-relative-only
B1 baseline; bootstrap 95\,\% CIs in brackets ($B = 1000$).
\vspace{-0.6em}
\end{table}

\begin{table}[!t]
\caption{Cross-session feature-group ablation verdict per fold (H0).}
\label{tab:cs-disambig}
\centering
\renewcommand{\arraystretch}{1.05}
\footnotesize
\begin{tabular}{l c c c c}
\toprule
Fold & B5 & B5$^\prime$ & B5$^{\prime\prime}$ & verdict \\
\midrule
F1 & 9.30 & 8.91 & 11.44 & LiDAR removable \\
F2 & 8.80 & 8.47 & 9.13 & LiDAR removable \\
F3 & 9.28 & 9.37 & 10.69 & LiDAR removable \\
\bottomrule
\end{tabular}
\\[3pt]
\raggedright
\footnotesize
Overall RMSE in dB; tolerance 0.5\,dB.
\vspace{-0.6em}
\end{table}

F3 nominally clears the $+1\,$dB threshold on the in-FOV stratum
($+1.14\dB$), but the feature-group ablation verdict is still
``LiDAR removable'' because B5$^\prime$ ties B5 within tolerance.
F3's apparent LiDAR boost is a substitution effect for AP-relative
features rather than a unique contribution: the LiDAR features are
encoding the same AP-direction information that B5$^\prime$ encodes
more compactly. This is exactly the distinction the feature-group
ablation procedure is designed to surface.

%-------------------------------------------------------------------- 5.3
\subsection{Hyperparameter robustness}
Table~\ref{tab:hyperparam} reports F1 in-FOV $\dlidar$ across six
configurations. Every configuration is negative; the range is
$[-1.90, -0.29]\dB$. The two random-validation variants run to their
early-stop limits (1{,}999 and 4{,}999 boosting rounds vs.\ 49 and 205
for their chronological-validation counterparts), which is a clear
over-fitting signature: random validation lets the model memorize
trajectory-correlated patterns that chronological validation correctly
excludes. The chronological protocol is therefore the more
conservative one, and the negative result is not an artifact of early
stopping.

\begin{table}[!t]
\caption{Hyperparameter robustness on F1 in-FOV.}
\label{tab:hyperparam}
\centering
\renewcommand{\arraystretch}{1.05}
\footnotesize
\begin{tabular}{l c c c c}
\toprule
Config & RMSE(B1) & RMSE(B5) & $\dlidar$ & clears +1\,dB? \\
\midrule
H0      & 8.01 & 8.53  & $-0.52$ & no \\
H1      & 7.86 & 8.15  & $-0.29$ & no \\
H2      & 8.15 & 8.59  & $-0.45$ & no \\
H3      & 8.58 & 9.19  & $-0.61$ & no \\
H1\_rv  & 7.86 & 9.76  & $-1.90$ & no \\
H2\_rv  & 8.15 & 9.38  & $-1.23$ & no \\
\bottomrule
\end{tabular}
\\[3pt]
\raggedright
\footnotesize
H0--H3 and the random-validation variants H1\_rv, H2\_rv are
defined in Section~\ref{sec:method-controls}. None clear the
$+1\,$dB threshold.
\vspace{-0.6em}
\end{table}

%-------------------------------------------------------------------- 5.4
\subsection{Framework cross-check (LightGBM)}
On F1 in-FOV, LightGBM with default hyperparameters (\texttt{num\_leaves}=31)
gives $\dlidar = -0.58\dB$; with depth-equivalent (\texttt{num\_leaves}=15)
$\dlidar = -0.93\dB$. Both are strictly more negative than any XGBoost
configuration. Across eight cells (six XGBoost configs $\cup$ two LightGBM
configs), $\dlidar$ is negative on every cell. The negative result is not
framework-specific.

%-------------------------------------------------------------------- 5.5
\subsection{Cross-session LiDAR-feature permutation control}\label{sec:cs-perm}
We retrain the full B5 model on F1 with the 19 ego-frame LiDAR
columns block-shuffled within each training session: each LiDAR
feature's marginal distribution and the joint distribution among LiDAR
features are preserved, but row-level alignment with telemetry,
position, AP-relative geometry, and the target is destroyed. If the
shuffled model achieves a $\dlidar$ comparable to the real-LiDAR fit,
the real-LiDAR $\dlidar$ was not driven by row-aligned LiDAR
information. Under both H0 and H1 hyperparameters,
$|\dlidar^{\text{real}} - \dlidar^{\text{shuffled}}| \le 0.30\dB$
(H0: $-0.52$ vs.\ $-0.37$; H1: $-0.29$ vs.\ $+0.01$), well within
the 0.5\,dB tolerance. The cross-session B5 model is not extracting
row-aligned LiDAR information beyond the marginal distributions of the
LiDAR features.

%==========================================================================
\section{Within-Session Experiment}\label{sec:within-session}
%==========================================================================

The within-session experiment evaluates whether LiDAR features add
value when raw position is available --- the regime in which a
deployed model would normally operate inside a fixed workspace.
Five-fold leave-region-out evaluation is performed on session S1
only, since S1 has by far the largest spatial coverage and sample
count of the three sessions. Regions are obtained by $K$-means
clustering on $(x_m, y_m)$ with $k = 5$ (Fig.~\ref{fig:regions});
the five resulting region splits give per-region row counts ranging
from 24{,}593 (R-4) to 80{,}365 (R-1). Each fold trains the
predictive models on four regions and evaluates on the held-out
region. Six ablation variants are evaluated, summarized in
Table~\ref{tab:wlro-variants}. The model's reliance on LiDAR is
quantified by the within-session ablation effect
$\dlidarw = \mathrm{RMSE}(\mathrm{W2}) - \mathrm{RMSE}(\mathrm{W4})$:
positive values indicate that LiDAR features improve over the
LiDAR-free W2 baseline.

\begin{table}[!t]
\caption{Within-session ablation variants.}
\label{tab:wlro-variants}
\centering
\renewcommand{\arraystretch}{1.05}
\setlength{\tabcolsep}{4pt}
\footnotesize
\begin{tabular}{l c l}
\toprule
Variant & \# feat. & Feature groups \\
\midrule
W0                  & 2  & position only \\
W1                  & 10 & W0 + telemetry \\
W2 ($\equiv$ W4$^\prime$) & 15 & W1 + AP-relative (no LiDAR) \\
W3                  & 20 & W2 + LiDAR scalar \\
W4                  & 34 & W3 + LiDAR sectoral (full model) \\
W4$^{\prime\prime}$ & 29 & W4 minus AP-relative \\
\bottomrule
\end{tabular}
\\[3pt]
\raggedright
\footnotesize
W2 is the LiDAR-free baseline; W4 is the full model.
W4$^{\prime\prime}$ removes AP-relative features and is the
counterpart of W2 in the feature-group ablation procedure of
Section~\ref{sec:method-disamb}.
\vspace{-0.6em}
\end{table}

% TODO: figure showing K-means regions on S1
\begin{figure}[!t]
\centering
% \includegraphics[width=\columnwidth]{figs/regions_15_03.pdf}
\fbox{\parbox[c][2.4cm][c]{0.92\columnwidth}{\centering\itshape
[Placeholder Fig.~\ref{fig:regions}: $K$-means $k{=}5$ partition of
the S1 trajectory in $(x_m, y_m)$. R-1 is closest to the AP; R-5 is
farthest.]}}
\caption{Within-session spatial regions on S1.}
\label{fig:regions}
\vspace{-1.0em}
\end{figure}

%-------------------------------------------------------------------- 6.2
\subsection{Headline results}
Table~\ref{tab:wlro-headline} reports per-fold in-FOV $\dlidarw$.
One of the five folds, R-4, nominally clears the $+1\,$dB
threshold ($+1.085\dB$); the other four are zero or negative. R-1
is severely negative ($-9.06\dB$ in-FOV): when the closest-to-AP
region is held out, the H0 (depth-6) model with LiDAR over-fits
sharply on the remaining regions and degrades on R-1. The
feature-group ablation verdict at the H0 configuration
(Table~\ref{tab:wlro-disambig}) is heterogeneous: R-3 and R-5 say
LiDAR removable; R-4 says AP-relative removable; R-1 and R-2 are
mixed. Without the robustness controls described below, the
within-session story would read as ``LiDAR helps in one of five
regions''; the controls invalidate that reading.

\begin{table}[!t]
\caption{Within-session in-FOV $\dlidarw$ per fold under H0.}
\label{tab:wlro-headline}
\centering
\renewcommand{\arraystretch}{1.05}
\footnotesize
\begin{tabular}{l c c c c}
\toprule
Fold & $n_{\text{test}}$ & RMSE(W2) & RMSE(W4) & $\dlidarw$ (95\,\% CI) \\
\midrule
R-1 & 57{,}270 & 9.51 & 18.57 & $-9.06$ ($[-9.10, -9.01]$) \\
R-2 & 15{,}516 & 5.19 & 6.08  & $-0.89$ ($[-0.92, -0.85]$) \\
R-3 & 30{,}974 & 4.39 & 4.52  & $-0.14$ ($[-0.16, -0.12]$) \\
R-4 & 10{,}687 & 3.11 & 2.02  & $\mathbf{+1.09}$ ($[+1.07, +1.10]$) \\
R-5 & 20{,}245 & 4.74 & 4.79  & $-0.04$ ($[-0.06, -0.03]$) \\
\bottomrule
\end{tabular}
\vspace{-0.6em}
\end{table}

\begin{table}[!t]
\caption{Within-session feature-group ablation verdict per fold (H0).}
\label{tab:wlro-disambig}
\centering
\renewcommand{\arraystretch}{1.05}
\footnotesize
\begin{tabular}{l c c c l}
\toprule
Fold & W4 & W4$^\prime$ & W4$^{\prime\prime}$ & verdict \\
\midrule
R-1 & 17.53 & 8.92 & 16.31 & mixed/unclear \\
R-2 & 6.06  & 5.49 & 6.67  & mixed/unclear \\
R-3 & 4.57  & 4.18 & 6.61  & LiDAR removable \\
R-4 & 4.69  & 5.76 & 5.00  & AP-rel.\ removable \\
R-5 & 10.68 & 10.60 & 11.40 & LiDAR removable \\
\bottomrule
\end{tabular}
\vspace{-0.6em}
\end{table}

%-------------------------------------------------------------------- 6.3
\subsection{Spatial autocorrelation in absolute RMSE}
Within-session RMSE values cannot be read as straightforward
generalization estimates: signal\_power is spatially autocorrelated,
and the held-out region rows are typically only a few metres away
from the nearest training rows. The S1 semivariogram has empirical
range $\approx 20$\,m, while the $K$-means regions span 5--13\,m on a
side, so training rows up to $\sim$10\,m from any held-out region row
are spatially correlated with the test set. To bound the magnitude of
this contamination we apply a buffer-zone test on R-1: dropping all
training rows within 1\,m of any held-out R-1 row inflates the W0 RMSE
by $+12\dB$, the W1/W2 RMSEs by $+1.5\dB$, and the W4 in-FOV by
$+2.3\dB$.

The contamination affects the LiDAR-free W2 baseline and the
LiDAR-augmented W4 model similarly within a fold, so the
\emph{contrast} $\dlidarw$ is not inflated by spatial autocorrelation
alone. More notably, the buffer-zone correction on R-4 \emph{shrinks}
its in-FOV $\dlidarw$ from $+1.09\dB$ to $-1.27\dB$: the contrast
moves in the direction unfavourable to the LiDAR model.

%-------------------------------------------------------------------- 6.4
\subsection{R-4 robustness diagnostic}
We test R-4's nominal $+1.09\dB$ in-FOV $\dlidarw$ under four
deployment-relevant corrections (Table~\ref{tab:r4-diag}). The buffer
correction alone gives $\dlidarw = -1.27\dB$, indicating that the
LiDAR model was relying on training rows in the spatial neighbourhood
of the held-out region to fit patterns that do not generalize when
those rows are removed. The depth-4 hyperparameter swap alone gives
$\dlidarw = -0.33\dB$, indicating that the H0 (depth-6) trees
over-fit LiDAR features which depth-4 trees correctly ignore. Joint
application of buffer and depth-4 gives $\dlidarw = -0.83\dB$: the two
corrections compound in the same direction rather than cancelling each
other out, ruling out the alternative interpretation that they
correct for opposing biases.

\begin{table}[!t]
\caption{R-4 four-correction robustness diagnostic. In-FOV $\dlidarw$
(dB). The R-4 nominal $+1.09\dB$ does not survive any correction; the
joint correction confirms the two individual corrections compound. The
permutation row reports the within-session-shuffled-LiDAR fit; its
matched real-LiDAR fit gives $\dlidarw = +1.11\dB$, so the gap between
the two is $0.98\dB$, exceeding the 0.5\,dB tolerance (see prose).}
\label{tab:r4-diag}
\centering
\renewcommand{\arraystretch}{1.1}
\setlength{\tabcolsep}{3pt}
\footnotesize
\begin{tabular}{l l r}
\toprule
Correction & probes & in-FOV $\dlidarw$ \\
\midrule
none (main run)            & ---                                  & $+1.09$ \\
1-m buffer                 & spatial autocorrelation              & $-1.27$ \\
depth-4 (H1)               & tree-depth over-fitting              & $-0.33$ \\
depth-4 + buffer           & both, simultaneously                 & $-0.83$ \\
LiDAR permutation          & row-aligned vs.\ marginal LiDAR      & $+0.13$ \\
\bottomrule
\end{tabular}
\vspace{-0.6em}
\end{table}

The within-session LiDAR-feature permutation control (last row of
Table~\ref{tab:r4-diag}) deserves a careful reading. We re-train W4 on
R-4 with the 19 ego-frame LiDAR columns block-shuffled (the
within-session counterpart of the cross-session control in
Section~\ref{sec:cs-perm}). The shuffled-feature
$\dlidarw = +0.13\dB$, while the matched real-LiDAR
$\dlidarw = +1.11\dB$. The 0.98\,dB gap exceeds the 0.5\,dB tolerance:
R-4's H0 no-buffer fit \emph{does} contain a small genuine
row-aligned LiDAR signal, and we report this signal explicitly rather
than pooling it into the negative result. However, this signal is
contingent on training data spatially adjacent to the test region (it
disappears under the 1-m buffer) and on tree depth $\ge 6$ (it
disappears under depth-4 trees). Under any of the three
deployment-relevant corrections --- buffer alone, depth-4 alone, or
their joint application --- $\dlidarw$ falls below the $+1\,$dB
threshold. The negative deployment-relevant claim therefore holds,
with explicit acknowledgement that R-4 contains a small genuine but
non-deployment-generalizing LiDAR signal.

%-------------------------------------------------------------------- 6.5
\subsection{Best-case within-session prediction accuracy}
The combination of corrections that produces the most defensible
deployment-relevant estimate --- depth-4 hyperparameters with R-4 as
the held-out region, in the in-FOV stratum --- gives the LiDAR-free
W2 model an RMSE of $1.99\dB$ on \texttt{signal\_power}. The 15
features of W2 are AGV position $(x_m, y_m)$, eight telemetry
features, and five AP-relative features. The full 34-feature W4 model
on the same fold and stratum has an RMSE of $2.33\dB$. Both numbers
sit approximately 3\,dB below the 0.5-m position-binning noise floor
of 4.95\,dB (Section~\ref{sec:noisefloor}), so the model's predictive
power is exclusively explained by sub-cell information --- position at
sensor resolution, telemetry, and AP-relative geometry --- rather
than by coarse spatial binning. LiDAR features cannot contribute
information that is not already encoded by the 15-feature stack.

%==========================================================================
\section{Feature Attribution Analysis}\label{sec:shap}
%==========================================================================

We complement the ablation-based analysis with a model-internal view
using TreeSHAP. The cross-session B5 fits and the within-session R-3
W4 fit are analysed; R-3 is chosen for the within-session view
because its overall RMSE of $4.57\dB$ is within $1\dB$ of the median
across the five within-session folds, so its attribution patterns are
representative rather than dominated by an outlier.

\subsection{Group importance}
The cumulative $|\mathrm{SHAP}|$ over the features in each group is
dominated by the AP-relative group on every fit
(Fig.~\ref{fig:shap-groups}). On the cross-session folds, AP-relative
contributes $5.6$--$8.0$ in cumulative $|\mathrm{SHAP}|$ (5 features;
per-feature $1.1$--$1.6$), whereas the LiDAR sectoral group
contributes $0.8$--$1.7$ across 14 features (per-feature
$0.06$--$0.12$). On the within-session R-3 W4 fit the contrast is
similar: AP-relative cumulative $|\mathrm{SHAP}|$ is $5.6$ over $5$
features against $1.5$ over $14$ for LiDAR sectoral --- a
ten-fold ratio in per-feature contribution. The ordering is
consistent with the ablation evidence in
Sections~\ref{sec:cross-session} and~\ref{sec:within-session}.

% TODO: figure showing SHAP feature group importance per fold
\begin{figure}[!t]
\centering
% \includegraphics[width=\columnwidth]{figs/shap_group_importance.pdf}
\fbox{\parbox[c][2.5cm][c]{0.92\columnwidth}{\centering\itshape
[Placeholder Fig.~\ref{fig:shap-groups}: Mean $|\mathrm{SHAP}|$ per
feature group, per cross-session fold.]}}
\caption{SHAP feature-group importance per cross-session fold.}
\label{fig:shap-groups}
\vspace{-0.8em}
\end{figure}

\subsection{Sign-of-effect consistency across folds}
We compute the Spearman correlation between each feature's value and
its SHAP value, per fold, and flag a feature as sign-consistent if
$|\rho| > 0.05$ with the same sign on all three cross-session folds.
Three of the four AP-relative features pass this test
(\texttt{dist\_to\_AP}, \texttt{sin\_AP},
\texttt{cos\_AP}); only \texttt{clutter\_frac\_toward\_AP}
flips sign across folds. Among the 19 LiDAR features, only four pass
(\texttt{mean\_dist\_mm}, \texttt{dist\_p90\_mm},
\texttt{mean\_dist\_sector\_6\_mm},
\texttt{mean\_dist\_sector\_7\_mm}). The asymmetry --- $75\%$ of
AP-relative features against $21\%$ of LiDAR features --- is
informative: AP-relative features encode a transferable physical
relationship that the gradient-boosted model picks up consistently
across folds, whereas the LiDAR-derived environmental geometry does
not, on this dataset.

\subsection{Spatial maps}
Plotting the per-cell dominant feature group in $(x_m, y_m)$ shows
that AP-relative and position together cover most of the AGV
trajectory on every fold; LiDAR sectoral never dominates a
contiguous spatial region. This is the per-cell counterpart of the
group-importance result above and provides a direct visual
confirmation.

%==========================================================================
\section{Discussion}\label{sec:discussion}
%==========================================================================

%-------------------------------------------------------------------- 7.1
\subsection{Why the negative result is informative}
Industrial WiFi propagation is governed by waveguiding through
corridors and constructive/destructive multipath at scales much larger
than typical LiDAR feature scales~\cite{cena2025wifi}. The negative
result is consistent with this picture: ego-frame LiDAR features over
the empirical 222$^{\circ}$ valid sector capture local geometry on
the metre scale, but the dominant propagation effects in the workspace
operate over distances substantially larger than the LiDAR's
sensitivity range and are not localized to the AGV's immediate
surroundings. AP-relative geometry derived from surveyed AP
coordinates, by contrast, encodes the global physical relationship
between AGV and AP that the gradient-boosted model can exploit
directly. The SHAP attribution analysis (Section~\ref{sec:shap})
supports this reading: AP-relative features dominate the model's
predictions across every fold, and the per-feature contribution of
LiDAR sectoral features is the smallest of any group.

%-------------------------------------------------------------------- 7.2
\subsection{Deployment guidance}
Operators with surveyed AP coordinates can predict
\texttt{signal\_power} at sub-2-dB within-session RMSE using a
15-feature stack: $(x_m, y_m)$ + 8 telemetry features + 5 AP-relative
features. The model is XGBoost with depth-4, $\eta=0.05$,
$\lambda=1.0$, $n_\text{est}=2000$, early-stop $= 100$; anomaly
handling is a single threshold $T^{\star} = 35$ on
\texttt{nns\_position\_confidence}. LiDAR-derived environment features
need not be in the stack: not for cross-session generalization (LiDAR
removable on every fold), and not for within-session generalization
under proper testing (no fold's nominal LiDAR boost survived a
deployment-relevant correction). The operator can therefore skip
LiDAR sensor-fusion engineering for link-quality prediction, the
sectoral feature pipeline, and in-FOV/out-of-FOV stratified models
with LiDAR-derived FOV. None added value over the 15-feature
AP-aware baseline in this deployment scenario.

%-------------------------------------------------------------------- 7.3
\subsection{Limitations}\label{sec:limitations}
Several limitations bound the scope of this paper's claims.
\textit{Single facility, single AP, three sessions.} The experiments
are conducted on one AGV in one industrial facility, across three
routes (one in a different sub-area); cross-facility replication is
left to future work. \textit{Single front LiDAR.} The AGV has a
single front-facing 270$^{\circ}$ LiDAR (1{,}350 active beams,
222$^{\circ}$ empirical valid sector), whereas production AGVs
typically carry front+rear LiDARs; the in-FOV stratum used
throughout is the dual-LiDAR-equivalent lower bound, and the negative
result holds in both strata. \textit{Spatial autocorrelation.} The
S1 semivariogram range ($\approx 20$\,m) is comparable to the
within-session region size and a 1\,m buffer does not fully
decorrelate, so the absolute within-session RMSE values within a
fold should be read as fold-conditional rather than as
deployment-grade generalization estimates --- though the
$\dlidarw$ contrast is not inflated by this alone
(Section~\ref{sec:within-session}-C). \textit{Hyperparameter
dependence.} The H0 depth-6 trees over-fit on some folds (notably
R-1, where depth-4 improves overall RMSE by 4.76\,dB); the cleanest
deployment-relevant numbers come from depth-4. The negative result
holds under all hyperparameter configurations tested but is more
decisive at depth-4. \textit{SHAP is descriptive, not causal.} The
SHAP analyses describe the trained models' learned relationships,
not underlying physical causes.
\textit{Single-threshold anomaly handling.} The histogram-saddle
threshold $T^{\star} = 35$ is applied uniformly across the three
sessions; per-session adaptive thresholds were not evaluated.

%==========================================================================
\section{Conclusion}\label{sec:conclusion}
%==========================================================================

This paper presented a use case study testing whether LiDAR-derived
environmental geometry can augment WiFi signal-strength prediction on
industrial AGVs beyond what AP-relative features achieve alone. The
study used cross-session and within-session generalization protocols
on 681{,}593 matched LiDAR--telemetry pairs from three data-gathering
sessions covering three routes. Across the experiments --- spanning
cross-session generalization, hyperparameter sensitivity, framework
sensitivity, within-session generalization with spatial and
over-fitting controls, and LiDAR-feature permutation controls ---
ego-frame LiDAR features did not improve on a 15-feature baseline of
position, telemetry, and AP-relative geometry derived from surveyed
AP coordinates. The F1-fold in-FOV ablation effect is
$-0.52\dB$ under the H0 configuration and lies in
$[-1.90, -0.29]\dB$ across six XGBoost configurations, with LightGBM
cross-checks confirming that the result is not framework-specific.
The cleanest within-session LiDAR-free model achieves an RMSE of
$1.99\dB$ on \texttt{signal\_power}, approximately 3\,dB below the
dataset's 0.5-m position-binning noise floor. The contributions are
(i) a controlled empirical study of LiDAR-augmented WiFi signal
strength prediction in this deployment scenario, (ii) a feature-group
ablation procedure, and (iii) practical deployment guidance for
industrial AGV operators: when AP coordinates are surveyed, ego-frame
LiDAR features were not required for sub-2-dB within-session
prediction in the studied scenario. Future work includes
cross-facility replication, dual-LiDAR coverage, and pose-uncertainty
modelling.

%==========================================================================
\section*{Acknowledgment}
%==========================================================================
% TODO: confirm grant numbers
Co-funded by the European Union HORIZON TMA MSCA Doctoral Networks /
HORIZON-MSCA-2023-DN-01 / project TUAI / grant agreement
N\textdegree~101168344. Views and opinions expressed are those of the
author(s) only and do not necessarily reflect those of the European
Union or European Research Executive Agency. Neither the European
Union nor the granting authority is responsible for them. Co-funded
by the Statutory Research funds of the Department of Distributed
Systems and Informatic Devices, Silesian University of Technology,
Gliwice, Poland (Grants BK-244/RAu8/2025 02/110/BK 25/1036).

%==========================================================================
\bibliographystyle{IEEEtran}
\bibliography{ref}

\end{document}
